\subsection{An ordinary order-five obstruction at a split \texorpdfstring{$A_6^{20}$}{A6-power} chief step}

\begin{statement}{D.36}{Proved}{Every nonidentity power of an ordinary peripheral order-five shadow is obstructed}
Put
\[
 \begin{split}
 m&=61478762617650840213814367698461663428608,\\
 h&=61478762617650840528585632300833961934848.
 \end{split}
 \tag{D.36.1}
\]
There are explicitly specified finite groups, a split chief extension,
and an actual source epimorphism
\[
 1\longrightarrow A_6^{20}\longrightarrow\mathcal E
       \xrightarrow{\pi}\mathcal B\longrightarrow1,
 \qquad q_{\mathcal E}:\widehat F\twoheadrightarrow\mathcal E,
 \tag{D.36.2}
\]
with
\[
 |\mathcal B|=5662310400\cdot360^{100m},
 \qquad
 |\mathcal E|=5662310400\cdot360^{100m+20}.
 \tag{D.36.3}
\]
An actual automorphism \(\gamma\) of \(\mathcal B\) has order five and
preserves the three labelled peripheral conjugacy classes of
\(q_{\mathcal B}=\pi q_{\mathcal E}\), with multiplier one.
For every \(k=1,2,3,4\), the ordinary fibre
\[
 \left\{f:\widehat F\longrightarrow\mathcal E:
       \pi f=\gamma^kq_{\mathcal B},\
       f(x_i)\in q_{\mathcal E}(x_i)^{\mathcal E}
       \ (i=1,2,3)\right\}
 \tag{D.36.4}
\]
has weak solutions and has no epimorphic solution.  All base
conjugators, including their full centralizer ambiguity, are allowed
in this assertion.

The maximal pro-five quotients of \(\mathcal B\) and \(\mathcal E\)
are trivial.  Taking their actual products with a prescribed finite
five-group source \(L\) retains an obstructed branch
\(\gamma^k\times\beta^k\), for each \(k=1,2,3,4\), and an admitted
branch \(1\times\beta^k\), over the same multiplier-one pro-five
shadow \(\beta^k\).  In particular, every \(5\)-adic unit power of
the specified order-five finite component remains obstructed.
These maps concern a specified finite component; the simultaneous
admitted branch is part of the asserted comparison with D.1.
\end{statement}

\begin{proof}
We give the finite groups and maps before the exhaustion argument.
The large integers in \textup{(D.36.1)} count explicitly defined
finite sets; selecting elements from these sets means taking their
lexicographically first generator tuples in the fixed enumeration
below.  Thus the construction makes no choice of an unevaluated
infinite lifting problem.

\paragraph{The small base with its actual order-five automorphism.}
Let \(H=A_5\) act on
\[
 M=\{(a_0,\ldots,a_4)\in\mathbf F_2^5:\textstyle\sum a_i=0\},
 \qquad
 b_i=e_i+e_4\quad(0\leq i<4).
 \tag{D.36.5}
\]
Use the actual permutations
\[
 u=(01234),\qquad v=(01342),\qquad
 w=(uv)^{-1}=(03412),
 \qquad C=\langle(012)\rangle.
 \tag{D.36.6}
\]
All three peripheral permutations have order five and \(u,v\)
generate the sixty-element group.  Put \(V=M^4\), with diagonal
\(H\)-action, and give \(V\rtimes H\) the multiplication
\((a,s)(b,t)=(a+s b,st)\).  The actual source map is
\[
 q_0^{\,\mathrm{mod}}:\widehat F\longrightarrow D_0=V\rtimes H,
 \qquad
 x_1\longmapsto((b_0,b_1,b_2,b_3),u),\quad
 x_2\longmapsto(0,v).
 \tag{D.36.7}
\]
It is onto.  One direct certificate follows every Cayley edge of
\(\langle u,v\rangle\): the resulting 61 actual source relators,
evaluated in \(V\), and their \(H\)-translates span all sixteen
coordinates.  Therefore the source image contains \(V\) and
surjects onto \(H\).  The saved relator words and sixteen-dimensional
span check prove the assertion without a generation assumption
about a split extension.  The same finite calculation verifies
that every nonzero vector of \(M\) has an orbit spanning \(M\),
and that \(\operatorname{End}_H(M)=\mathbf F_2\).
Its crossed-homomorphism constraint matrix has rank four on the
eight generator coordinates: \(Z^1(H,M)=B^1(H,M)\) has dimension
four.  These give a second, compatible verification of
\textup{(D.36.7)}.

Define the multiplicity operator
\[
 T(a_0,a_1,a_2,a_3)
       =(a_3,a_0+a_3,a_1+a_3,a_2+a_3).
 \tag{D.36.8}
\]
Its four columns are \(2,4,8,15\) in binary coordinates.
It commutes with the actual \(H\)-action, \(T^5=1\), and \(T\ne1\).
Let \(G=\operatorname{Aut}(A_6)\), \(S=\operatorname{Inn}(A_6)\),
and \(O=G/S=\mathbf F_2^2\), with the actual multiplication and
outer coordinates of D.31.  Take as \(q_Q:\widehat F\twoheadrightarrow G\)
the map \(H_0\) of \textup{(D.32.4)--(D.32.5)}.  Its peripheral
indices in the complete generator-image enumeration of D.31
are \(0,24,1380\), and its outer values are \(1,2,3\).
Its complete peripheral centralizer images are
\[
 \overline{C_G(q_Q(x_1))}=\{0,1\},\quad
 \overline{C_G(q_Q(x_2))}=\{0,2\},\quad
 \overline{C_G(q_Q(x_3))}=\{0,3\}.
 \tag{D.36.9}
\]
The certificate enumerates all members of these three centralizers.

Set
\[
 D=D_0\times G,\qquad
 q_0=(q_0^{\,\mathrm{mod}},q_Q),\qquad
 \alpha(a,s,g)=(Ta,s,g),\qquad d_i=q_0(x_i).
 \tag{D.36.10}
\]
The map \(q_0\) is onto.  Indeed, \(D_0\) is perfect:
\([V,H]=V\) by the displayed orbit-span calculation, and \(A_5\)
is perfect.  The group \(G\) has no nontrivial perfect quotient.
For if \(N\lhd G\), simplicity of \(S\) gives either \(S\subset N\)
or \(N\cap S=1\); in the latter case \([N,S]=1\), whence
\(N\subset C_G(S)=1\).  Its proper quotients therefore factor
through \(O\), and \(G\) itself has nontrivial abelian quotient
\(O\).  Goursat's actual common-quotient map for the two onto
projections of \(q_0\) must consequently be trivial.
This proves the claimed product generation.
In particular \(|D|=60\cdot2^{16}\cdot1440=5662310400\).

Write \(q_0^{\,\mathrm{mod}}(x_i)=(v_i,h_i)\).
The operator \(1-h_i\) is invertible on each copy of \(M\).
For \(j=0,\ldots,4\), define the literal module element
\[
 c_i^{(j)}=(1-h_i)^{-1}(T^j-1)v_i\in V\subset D.
 \tag{D.36.11}
\]
It satisfies
\[
 \alpha^j(d_i)=c_i^{(j)}d_i(c_i^{(j)})^{-1},\qquad
 c_i^{(0)}=c_i^{(5)}=0,\qquad
 c_i^{(j+1)}-c_i^{(j)}=\alpha^j(c_i^{(1)}).
 \tag{D.36.12}
\]
These identities use the displayed multiplication in \(V\rtimes H\);
they are also checked on every saved vector.  Thus \(\alpha\)
is an ordinary peripheral automorphism of exact order five.

\paragraph{A specified twenty-factor chief extension.}
For a vector in the basis \(b_0,\ldots,b_3\), put
\(\ell(a)=a_0+a_1+a_2\) and \(r(a)=a_3\).
The homomorphism
\[
 \epsilon:D_*=(V\rtimes C)\times G\longrightarrow O,\qquad
 \epsilon((a_0,a_1,a_2,a_3),s,g)
       =(\ell(a_0),r(a_1))+\overline g
 \tag{D.36.13}
\]
is well defined: the two module functionals are \(C\)-invariant.
They are encoded by the sixteen-bit masks \(7,128\).
Let \(\Omega=D_*\backslash D=C\backslash H\), of size twenty,
and choose representatives \(t_\omega=(0,s_\omega,1)\).
Here \(s_\omega\) is given by the signed-edge breadth-first tree
in the \(u,v\) Schreier graph, with edge order
\(x_1,x_2,x_1^{-1},x_2^{-1}\).  All group permutations, cosets,
and tree words are retained in the certificate.
Put
\[
 a(\omega,d)=\epsilon(t_\omega d t_{\omega d}^{-1}),\qquad
 E_*=\{(z,d)\in G^\Omega\rtimes D:
                   \overline{z(\omega)}=a(\omega,d)\}.
 \tag{D.36.14}
\]
The multiplication is literally
\[
 (z,d)(z',d')
       =\bigl((z(\omega)z'(\omega d))_\omega,dd'\bigr).
 \tag{D.36.15}
\]
The cocycle identity for \(a\) follows by cancelling the intermediate
representative.  Projection \(E_*\to D\) is onto: in each coordinate
choose an actual \(G\)-representative of the specified outer value.
Its kernel is \(A=S^\Omega\), and \(D\) permutes the twenty factors
transitively.  A normal subgroup of \(E_*\) contained in \(A\)
is a product of simple direct factors and is therefore \(1\) or
\(A\).  Thus this is an actual chief extension.

Put \(K=\ker q_0\) and
\(\mathcal H=q_0^{-1}(D_*)\).  Since \(\alpha\) preserves \(D_*\),
these two source subgroups are the same for all five maps
\(q_j=\alpha^jq_0\).  The index-twenty cover of the three-punctured
sphere has four cycles of length five over each peripheral.
It has twelve boundaries, genus five, and a free Schreier basis
of rank 21.  The twenty tree words will be denoted \(T_\omega\).
They are actual words in \(x_1,x_2\), whereas \(t_\omega\)
in \textup{(D.36.14)} is the chosen element of \(D\).
This difference is retained below.

For a homomorphism \(\varphi:\mathcal H\to G\) with
\(\overline\varphi=\epsilon q_j|_{\mathcal H}\), choose actual
\(b_{\omega,j}\in G\) satisfying
\[
 b_{*,j}=1,\qquad
 \overline b_{\omega,j}
       =\epsilon(q_j(T_\omega)t_\omega^{-1}).
 \tag{D.36.16}
\]
For example, take the least indexed \(G\)-element of that outer
value away from the identity coset.  The actual induced source
map \(p_\varphi:\widehat F\to E_*\), over \(q_j\), is
\[
 z_\varphi(g)(\omega)
   =b_{\omega,j}^{-1}
      \varphi(T_\omega gT_{\omega g}^{-1})
      b_{\omega g,j}.
 \tag{D.36.17}
\]
Multiplying two right sides cancels the intermediate \(b\) and
gives \textup{(D.36.15)}.  Taking outer coordinates gives
\(a(\omega,q_j(g))\), by \textup{(D.36.16)}.
Conversely restriction to the identity coordinate on
\(\mathcal H\), with the actual path values \(z(T_\omega)(*)\),
recovers this formula.  This is the needed nonabelian Shapiro
comparison, including its forced outer path coordinates.

\paragraph{The fixed peripheral classes and their full based transport.}
Let
\[
 \mathcal C_1=X^S,\qquad
 \mathcal C_2=q_Q(x_2)^S,\qquad
 \mathcal C_3=q_Q(x_3)^S,
 \tag{D.36.18}
\]
where \(X\) is the actual involution of \textup{(D.31.6)}.
Their indices in the relative class table are \(12,7,5\);
their sizes are \(15,90,90\), and their outer values are \(1,2,3\).
The action of \(O\) on \(S\)-classes means conjugation by any
actual representative in \(G\); it is independent of that
representative.

Choose the least coset in each peripheral cycle as its start.
Write \(T_\nu\) for its tree word and
\[
 q_0^{\,\mathrm{mod}}(T_\nu)=(v_\nu,s_\nu),\qquad
 \delta_\nu=T_\nu x_i^5T_\nu^{-1}\in\mathcal H
 \quad(\nu\text{ lies over }x_i).
 \tag{D.36.19}
\]
There are four such \(\nu\) for each \(i\).  Define
\(\eta_\nu=2\) at the first \(x_1\)-cycle, which contains
the identity coset, and \(\eta_\nu=0\) elsewhere.
Specify the cycle-product class of a reference element
\(Z_i\in E_*\), over \(d_i\), to be
\[
 \bigl(\epsilon_V(v_\nu)+\eta_\nu\bigr)\cdot\mathcal C_i,
 \qquad
 \epsilon_V(a)=(\ell(a_0),r(a_1)).
 \tag{D.36.20}
\]
These actual reference elements exist coordinate by coordinate.
Along a length-five cycle choose the first four coordinate
elements with their required outer values, and choose the fifth
so that the ordered product is the least element of the class
in \textup{(D.36.20)}.  The required outer product is \(i\in O\),
so the last coordinate has exactly the required outer value.
The certificate gives all sixty coordinate elements.

For a source over \(q_j\) whose \(i\)-th peripheral is conjugate
to \(Z_i\) by an element over \(c_i^{(j)}\), the precise class
of \(\varphi(\delta_\nu)\) is
\[
 \mathcal T_{\nu,j}
  =
 \left(
   \overline{q_Q(T_\nu)}
   +\epsilon_V((T^j-1)v_\nu)
   +\epsilon_V(s_\nu c_i^{(j)})
   +\eta_\nu
 \right)\cdot\mathcal C_i.
 \tag{D.36.21}
\]
Here \(s_\nu c_i^{(j)}\) is the specified module action.
To verify the formula, the cycle product changes under a lift
of \(c_i^{(j)}\) by the outer value
\(\epsilon_V(s_\nu c_i^{(j)})\).  The based boundary value is
that cycle product conjugated by \(b_{\nu,j}\) from
\textup{(D.36.16)}, whose outer value is
\(\epsilon_V(T^jv_\nu)+\overline{q_Q(T_\nu)}\).
Adding the initial adjustment \(\epsilon_V(v_\nu)+\eta_\nu\)
in \textup{(D.36.20)} proves \textup{(D.36.21)}.
In particular the module term
\(\epsilon_V((T^j-1)v_\nu)\) is required even though all five
source maps have the same permutation action on the cosets.

For completeness the conjugacy criterion just used has a literal
inverse.  For two arrays over the same \(d_i\), let \(P,P'\)
be their ordered products around a cycle.  Choose \(a_0\in S\)
with \(P'=a_0Pa_0^{-1}\).  If the old and new coordinates are
\(z_\omega,z'_\omega\), successively put
\[
                  a_{\omega d_i}
                         =(z'_\omega)^{-1}a_\omega z_\omega.
 \tag{D.36.22}
\]
All \(a_\omega\) lie in \(S\), and the return to the starting
coordinate is \(a_0\) exactly because of the product equation.
This produces the actual kernel conjugator.  Thus
\textup{(D.36.21)} is both necessary and sufficient for the
specified local conjugator class.  The five saved source maps
verify this recurrence on every coordinate.

\paragraph{The complete finite count, with proper images removed.}
The twelve boundary words in the rank-21 Schreier basis are
literal words.  Across them every Schreier letter occurs once
positively and once negatively.  Introduce twelve boundary
generators \(\delta_\nu\).  Merging faces along eleven tree
edges gives a single quadratic relation with ten paired
letters and twelve boundary letters.  All these eliminations
are actual Tietze maps: from relations \(aU=1\) and
\(a^{-1}V=1\), substitute \(a=U^{-1}\) and retain \(UV=1\).
The remaining handles are extracted by the exact identity
\[
 \begin{gathered}
 R=aUbVa^{-1}Wb^{-1}Z,\qquad C=WVU,\\
 A=C^{-1}aU,\qquad B=bVU,\qquad
 C^{-1}RC=[A,B]ZWVU.
 \end{gathered}
 \tag{D.36.23}
\]
The substitutions are invertible:
\(a=CAU^{-1}\) and \(b=BU^{-1}V^{-1}\).
Cyclic changes of the remaining polygon conjugate previously
extracted handles explicitly and leave each boundary generator
unchanged.  The certificate checks every free-word identity,
the resulting relation, and the determinant \(-1\) of the
21-element integral basis comparison.  It also inverts these
maps in \(G\), producing actual homomorphisms on the original
Schreier generators.

For each \(j\), let \(\mathscr X_j\) consist of the onto maps
\(\varphi:\mathcal H\to G\) with the exact outer character
\(\epsilon q_j\) and all twelve conditions
\(\varphi(\delta_\nu)\in\mathcal T_{\nu,j}\).
The complete count is
\[
 \begin{array}{c|ccccc}
 j&0&1&2&3&4\\\hline
 |\mathscr X_j|&360h&360m&360m&360m&360m .
 \end{array}
 \tag{D.36.24}
\]
We describe the entire finite calculation so that generation
and framed transport are part of its proof.

First enumerate every subgroup \(L\leq G\) mapping onto \(O\).
Such a subgroup contains a pair with outer values \(1,2\).
The first member can be moved, by \(S\)-conjugation, into one
of its five inner conjugacy-class representatives.  The 1800
resulting pairs give all starting subgroups up to conjugacy.
Adjoining every inner element to each new subgroup exhausts
all possibilities: a subgroup surjecting onto \(O\) is
generated by such a starting pair and its intersection with
\(S\).  This gives ten conjugacy classes and exactly 452
actual subgroups, with order histogram
\[
 \begin{array}{c|rrrrrr}
 |L|&8&16&32&40&144&1440\\\hline
 \#L&180&180&45&36&10&1.
 \end{array}
 \tag{D.36.25}
\]
For these subgroups \(G\)-conjugacy and \(S\)-conjugacy coincide:
if \(L\) surjects onto \(O\) and \(g\in G\), choose
\(l\in L\) with \(\bar l=\bar g\); then
\(gLg^{-1}=(gl^{-1})L(gl^{-1})^{-1}\) and \(gl^{-1}\in S\).
Thus all prescribed boundary sets are preserved in using the
ten representatives.

For each \(L\), put \(L_0=L\cap S\).  In
\(\mathbf Z[L]^{L_0}\) compute the actual elements
\[
 K^L_{a,b}
   =\sum_{\substack{u,v\in L\\\bar u=a,\ \bar v=b}}[u,v],
 \qquad
 B^L_{\mathcal T}=\sum_{t\in\mathcal T\cap L}t.
 \tag{D.36.26}
\]
The handle outer values are those obtained from the literal
words in \textup{(D.36.23)}, and the boundary sets are exactly
\textup{(D.36.21)}.  The coefficient of \(1\) in the ordered
product of the five appropriate \(K^L_{a,b}\) and the twelve
inverse-boundary sums counts all such maps into \(L\).
This is just expansion of the actual surface relation; it
uses no inference from a rational character formula.
The relative class-algebra structure constants are computed
by counting \(x\) in one class with \(x^{-1}g\) in the other.
Every coefficient is an integer.

In increasing subgroup order subtract the counts of maps
whose image is each actual proper subgroup:
\[
 E_L=N_L-\sum_{J<L,\ \bar J=O}E_J.
 \tag{D.36.27}
\]
Because the prescribed outer character is onto \(O\), this
list includes every proper image.  The count for \(L=G\)
is \textup{(D.36.24)}.  The positive excess of kernel classes is
\[
                         h-m=314771264602372298506240.
 \tag{D.36.28}
\]
As a separately executed check, omitting the one fixed flip
\(\eta_\nu=2\) gives \(360(m,h,h,h,h)\).
An earlier exploratory calculation omitted the module-path
term in \textup{(D.36.21)}; its programs and outputs are
preserved with the suffix \texttt{incomplete\_framing}.
They are not used in \textup{(D.36.24)}.
The corrected calculation includes both the term and the
two opposite patterns.

\paragraph{The exact comparison with source kernels.}
Every \(\varphi\in\mathscr X_j\) sends \(K\) onto \(S\).
Its image on \(K\) is normal in \(G\), lies in \(S\), and
is therefore \(1\) or \(S\).  If it were \(1\), \(\varphi\)
would factor through
\[
                 \mathcal H/K=D_*=(V\rtimes C)\times G.
 \tag{D.36.29}
\]
The first factor is solvable.  Its image would be a solvable
normal subgroup of \(G\), hence trivial: a normal subgroup
meeting \(S\) nontrivially contains the nonsolvable simple
group \(S\), while a normal subgroup disjoint from \(S\)
centralizes \(S\) and is trivial.  The factored map would
then kill \(V\), contradicting its prescribed nonzero
outer character \(\epsilon_VT^j\) on \(V\).
This proves \(\varphi(K)=S\).

For the five ordered module-character pairs, the exact
stabilizer and transporter calculation is
\[
 \{s\in H:\epsilon_VT^j s=\epsilon_VT^k
                    \text{ on }V\}
   =\begin{cases}C,&j=k,\\\varnothing,&j\ne k.\end{cases}
 \tag{D.36.30}
\]
It is verified on all sixteen basis vectors for all sixty
actual permutations and all 25 ordered pairs.  Equivalently,
the two masks \(7,128\) on the specified module copies and the
matrix \textup{(D.36.8)} give the displayed finite calculation.

Let \(R_\varphi=\ker(\varphi|_K)\).  Its normalizer in
\(\widehat F\) is exactly \(\mathcal H\).
Indeed, the normalizer contains \(K\), and its conjugation
action on \(K/R_\varphi=S\) gives an actual map to
\(\operatorname{Aut}(S)=G\) extending \(\varphi\).
Taking outer classes gives an extension of
\(\epsilon q_j|_{\mathcal H}\).
Every subgroup of \(D\) containing \(D_*\) is the preimage
of a subgroup of \(H\) containing \(C\), since \(D_*\)
already contains \(V\times G\).
Its action must fix the character pair on \(V\);
\textup{(D.36.30)} forces that subgroup to be \(C\).
Thus the normalizer equals \(\mathcal H\).
In particular the twenty coordinate kernels in
\textup{(D.36.17)} are distinct.

If two restrictions to \(K\) have equal kernels, their
onto maps \(K\to S\) differ by an automorphism of \(S\),
that is, by an actual element \(g\in G\).
The two extended maps on \(\mathcal H\) then differ by
conjugation by \(g\): their conjugation actions agree on
the onto subgroup \(S\), and \(C_G(S)=1\).
Conversely such conjugation preserves the kernel.
Within \(\mathscr X_j\), exactly 360 elements of this
\(G\)-orbit remain in the twelve prescribed classes.
Indeed the common stabilizer in \(O\) of the three classes
in \textup{(D.36.18)} is \(0\); the independent class
transports in \textup{(D.36.21)} do not change that
stabilizer because \(O\) is abelian.
The \(S\)-action on onto maps is free.
Thus \(\mathscr X_j\) represents exactly
\[
                           h,m,m,m,m
 \tag{D.36.31}
\]
different kernels on \(K\).

These are also distinct \(\widehat F\)-orbits of kernels.
Any conjugator between two such kernels must normalize
\(\mathcal H\), by the normalizer equality just proved.
Its image in \(H\) belongs to \(N_H(C)\).
The \(V\)-part acts trivially on the \(C\)-invariant
characters, and the \(G\)-part acts trivially on \(O\).
Consequently the required character comparison is precisely
\textup{(D.36.30)}.  For equal colours it forces the
conjugator into \(\mathcal H\), where each kernel is already
normal; for different colours no conjugator exists.

The simple-product generation argument is exact.
An onto map to each of finitely many copies of a nonabelian
simple group, with pairwise distinct kernels, maps onto
their product.  Inductively, a proper subdirect subgroup
of \(S^{n-1}\times S\) would be the graph of an onto map
\(S^{n-1}\to S\).  The images of the direct factors are
commuting normal subgroups of \(S\); simplicity and
centerlessness force this map to factor through one
coordinate.  Its kernel would coincide with an earlier
one.  Applying this argument to the twenty coordinate
kernels proves that every map in \(\mathscr X_j\) induces
an epimorphism in \textup{(D.36.17)}, and applying it
across colours proves all product-generation assertions
used below.

\paragraph{The entire ordinary centralizer union.}
This step is independent of the numerical count.
Since \(1-h_i\) is invertible on \(V\), the element
\((v_i,h_i)\) is conjugate by a unique module translation
to \((0,h_i)\).  Its centralizer is conjugate to
\(C_H(h_i)\), which is the cyclic group of order five.
Hence
\[
 C_D(d_i)=\langle q_0^{\,\mathrm{mod}}(x_i)\rangle
                       \times C_G(q_Q(x_i)).
 \tag{D.36.32}
\]
This group is generated by \(d_i\) and the elements
\((1,c)\) with \(c\in C_G(q_Q(x_i))\).
Conjugation by the actual lift \(Z_i\) of \(d_i\)
fixes \(Z_i\).  For \((1,c)\), the permutation of
\(\Omega\) is the identity and the outer value in
every coordinate is \(\bar c\).  By
\textup{(D.36.9)}, this is \(0\) or \(i\).
Every \(S\)-class in outer coset \(i\) is fixed by
the outer action of \(i\): use a representative of
that very class as the conjugating element.
Thus conjugation by any lift of \((1,c)\) preserves
each cycle-product \(S\)-class of \(Z_i\).
The explicit recurrence \textup{(D.36.22)} shows
that it preserves its full \(A\)-conjugacy class.
Changing the chosen lift only adds an \(A\)-conjugation.
It follows that every element of the complete base
centralizer \textup{(D.36.32)} preserves the
\(A\)-class of \(Z_i\).

Consequently, for every \(j\), allowing all ordinary
base conjugators from \(d_i\) to \(\alpha^j(d_i)\)
gives exactly the class condition
\textup{(D.36.21)}.  Multiplying a specified
\(c_i^{(j)}\) by any centralizer element cannot
add another cycle-product profile.  This proves
the full ordinary comparison used by the exhaustion
argument, rather than only a comparison for one
chosen tuple of conjugators.

\paragraph{The five-colour finite product and its source.}
For \(j\in\mathbf Z/5\mathbf Z\), let
\[
 E_j=E_*\times_{D,\alpha^j}D
       =\{(e,d):\operatorname{pr}(e)=\alpha^j(d)\}.
 \tag{D.36.33}
\]
Order the onto maps in \(\mathscr X_j\) by their
21 generator values in the actual enumeration of \(G\).
In each equivalence class under \(G\)-conjugation
retain its least member still satisfying
\textup{(D.36.21)}.  The preceding kernel comparison
shows that these lists have lengths \(h,m,m,m,m\)
and have disjoint kernel orbits across colours.
Use the first \(m\) members for colour zero and
all \(m\) members for each other colour.
Denote the induced source maps by
\(p_{j,r}:\widehat F\to E_*\), over \(q_j\),
for \(1\leq r\leq m\).
Use the next colour-zero member for a map \(p_*\).
This is a completely specified finite procedure
on tuples in \(G^{21}\); its required positive
excess is \textup{(D.36.28)}.

Define the actual fibre product
\[
 \mathcal B
  =\left\{((e_{j,r}),d):
       \operatorname{pr}(e_{j,r})=\alpha^j(d)
       \text{ for every }j,r\right\},
 \qquad
 q_{\mathcal B}(g)=((p_{j,r}(g)),q_0(g)).
 \tag{D.36.34}
\]
The comparison of all coordinate kernels proves
that \(q_{\mathcal B}\) is onto.
Now define
\[
 \begin{split}
 \mathcal E&=\mathcal B\times_D E_*,\\
 q_{\mathcal E}(g)&=(q_{\mathcal B}(g),p_*(g)),\\
 \gamma((e_{j,r}),d)&=((e_{j+1,r}),\alpha(d)).
 \end{split}
 \tag{D.36.35}
\]
The same distinct-kernel argument proves that
\(q_{\mathcal E}\) is onto.  The shift in
\textup{(D.36.35)} has the required type because
\(\operatorname{pr}(e_{j+1,r})=\alpha^j(\alpha(d))\),
including \(j=4\).  It is an automorphism of
exact order five.
Forgetting the last \(E_*\)-coordinate gives
\textup{(D.36.2)}, with kernel \(S^{20}\).
It splits by the actual map taking that coordinate
to \(e_{0,1}\).  Its factor action is transitive,
so the kernel is chief, as for \textup{(D.36.14)}.
Counting the \(5m\) kernels of size \(360^{20}\)
proves \textup{(D.36.3)}.

The automorphism \(\gamma\) is ordinary peripheral.
For each \(i,j,r\), choose an actual conjugator
from \(Z_i\) to \(p_{j,r}(x_i)\) over \(c_i^{(j)}\).
These exist by \textup{(D.36.21)--(D.36.22)}.
The quotient of two successive such conjugators
lies over
\(c_i^{(j+1)}-c_i^{(j)}=\alpha^j(c_i^{(1)})\).
These \(5m\) coordinates therefore assemble into
one actual element of \(\mathcal B\) over
\(c_i^{(1)}\) conjugating \(q_{\mathcal B}(x_i)\)
to \(\gamma q_{\mathcal B}(x_i)\).
The cyclic closing coordinate uses \(c_i^{(5)}=0\)
and gives the same equality.  This establishes all
three peripheral conjugacies in the actual product.

\paragraph{Exhaustion for every nonidentity power.}
Suppose \(f:\widehat F\to\mathcal E\) belongs to
\textup{(D.36.4)} and is onto.  Its extra
\(E_*\)-coordinate \(e\) is an epimorphism over
\(q_k=\alpha^kq_0\).  Project an arbitrary
peripheral conjugator in \(\mathcal E\) to
\(D\).  Its value belongs to
\(c_i^{(k)}C_D(d_i)\).
The full centralizer comparison
\textup{(D.36.32)} proves that the corresponding
map on \(\mathcal H\) belongs to
\(\mathscr X_k\).  Its \(K\)-kernel orbit
therefore occurs among the \(m\) already selected
colour-\(k\) maps.

In the base map \(\gamma^kq_{\mathcal B}\),
these same \(m\) maps occur as the \(j=0\)
coordinates: the shift has put \(p_{k,r}\)
in those positions.  Thus some simple-coordinate
kernel of the extra \(E_*\)-map coincides with
a simple-coordinate kernel already present
in the base.  The image of \(K\) in the
kernel of \(\mathcal E\to D\) is contained in
the corresponding proper diagonal graph.
An epimorphism \(f\) would send \(K\) onto
that whole kernel, a contradiction.
This proves the emptiness of the generating
fibre for all four values of \(k\).

Weak lifts nevertheless exist.  For any
\(r\), compose \(\gamma^kq_{\mathcal B}\)
with the split section that uses \(e_{0,r}\)
as the extra coordinate.  The extra source
map is then \(p_{k,r}\).  Its peripheral
classes agree with those of \(p_*\) after
conjugators over \(c_i^{(k)}\); these may
be paired with the already constructed
base conjugators, giving actual elements
of \(\mathcal E\).
This supplies an element of
\textup{(D.36.4)} with complementary image.

\paragraph{The exact pro-five comparison and reproducible inputs.}
Every homomorphism from \(A_6\) or \(A_5\)
to a finite five-group is trivial, and a
homomorphism from the module \(V\) to a
five-group is trivial.  The quotient
\(G/S=O\) has order four.  Applying these
facts to the displayed normal sequences
shows directly that \(D,E_*,\mathcal B\),
and \(\mathcal E\) all have trivial maximal
pro-five quotient.
For a finite five-group source \(L\),
the maps \((q_{\mathcal E},q_L)\) and
\((q_{\mathcal B},q_L)\) are onto the
actual products: Goursat's common quotient
would be a five-group quotient of the first
factor and is therefore trivial.
Their maximal pro-five quotients are
exactly \(L\).  Projection to
\textup{(D.36.4)} obstructs every branch
\(\gamma^k\times\beta^k\); the actual
automorphism \(1_{\mathcal E}\times\beta^k\)
gives the admitted branch
\(1_{\mathcal B}\times\beta^k\).
This proves the stated comparison without
discarding the specified finite component.

The finite dependencies are the actual
\(A_6\) model and source table from D.31--D.32,
followed by the programs in the \texttt{code} directory
with common prefix \texttt{a\_peripheral\_} and suffixes
\texttt{a5\_deleted.py},
\texttt{a6\_outer\_subgroups.cpp},
\texttt{a6\_cup\_epi.py},
\texttt{surface20.py},
\texttt{a6\_boundary\_search.py}, and
\texttt{order5\_certificate.py}.
The authoritative surface data have suffix
\texttt{surface20\_framed}; the authoritative
count data have suffix
\texttt{boundary\_framed\_flip}.
The last certificate constructs one actual
onto surface map in each colour, inverts
the Tietze maps to its 21 original generator
values, computes the two induced
twenty-coordinate source arrays using
\textup{(D.36.17)}, and checks all twelve
cycle products and all sixty peripheral
conjugator coordinates for every colour.
It also checks all 25 transporter sets in
\textup{(D.36.30)} and the complete base
centralizer images.  The five-colour
exhaustion uses the full counts
\textup{(D.36.24)}, with the exact finite
lexicographic selection in
\textup{(D.36.34)}, rather than claiming
to enumerate the enormous product group.
\end{proof}
